\documentclass[11pt, letterpaper]{article}
% ============================================================
% preamble.tex — Shared LaTeX preamble for Learning to Autolens
% ============================================================
% Contains: packages, physics notation macros, formatting.
% Include in each module via % ============================================================
% preamble.tex — Shared LaTeX preamble for Learning to Autolens
% ============================================================
% Contains: packages, physics notation macros, formatting.
% Include in each module via % ============================================================
% preamble.tex — Shared LaTeX preamble for Learning to Autolens
% ============================================================
% Contains: packages, physics notation macros, formatting.
% Include in each module via \input{../preamble}
% ============================================================

% --- Packages ---
\usepackage[utf8]{inputenc}
\usepackage[T1]{fontenc}
\usepackage{amsmath, amssymb, amsthm}
\usepackage{physics}       % \vb, \grad, \div, \curl, \dd, etc.
\usepackage{mathtools}     % \coloneqq, \prescript
\usepackage{bm}            % \bm for bold math
\usepackage{graphicx}
\usepackage{float}
\usepackage{booktabs}      % Professional tables
\usepackage{hyperref}
\usepackage{cleveref}      % \cref for smart cross-refs
\usepackage{xcolor}
\usepackage[margin=1in]{geometry}
\usepackage{enumitem}      % Custom list formatting
\usepackage{fancyhdr}      % Headers and footers
\usepackage{tcolorbox}     % Colored boxes for key results
\usepackage{listings}      % Code listings

% --- Color scheme ---
\definecolor{codeblue}{RGB}{31, 119, 180}
\definecolor{codegray}{RGB}{128, 128, 128}
\definecolor{keyresult}{RGB}{39, 123, 69}
\definecolor{exercise}{RGB}{178, 102, 0}

% --- tcolorbox environments ---
\newtcolorbox{keyresult}[1][]{
  colback=keyresult!5, colframe=keyresult!70,
  fonttitle=\bfseries, title={Key Result}, #1
}

\newtcolorbox{pythonbox}[1][]{
  colback=codeblue!5, colframe=codeblue!50,
  fonttitle=\bfseries, title={PyAutoLens Implementation}, #1
}

\newtcolorbox{exercisebox}[1][]{
  colback=exercise!5, colframe=exercise!50,
  fonttitle=\bfseries, title={Exercise}, #1
}

% --- Code listing style ---
\lstset{
  language=Python,
  basicstyle=\ttfamily\small,
  keywordstyle=\color{codeblue}\bfseries,
  commentstyle=\color{codegray}\itshape,
  stringstyle=\color{keyresult},
  breaklines=true,
  frame=single,
  framerule=0.5pt,
  rulecolor=\color{codegray!50},
}

% --- Physics macros ---
% Vectors (bold upright)
\newcommand{\vtheta}{\bm{\theta}}       % Image-plane position
\newcommand{\vbeta}{\bm{\beta}}         % Source-plane position
\newcommand{\valpha}{\bm{\alpha}}       % Deflection angle
\newcommand{\vx}{\bm{x}}               % Generic position

% Lensing quantities
\newcommand{\thetaE}{\theta_{\rm E}}    % Einstein radius
\newcommand{\SigmaCr}{\Sigma_{\rm cr}}  % Critical surface density
\newcommand{\Dd}{D_{\rm d}}             % Angular diameter distance to lens
\newcommand{\Ds}{D_{\rm s}}             % Angular diameter distance to source
\newcommand{\Dds}{D_{\rm ds}}           % Angular diameter distance lens-source

% Statistical quantities
\newcommand{\chisq}{\chi^2}             % Chi-squared
\newcommand{\chisqred}{\chi^2_{\rm red}} % Reduced chi-squared
\newcommand{\Like}{\mathcal{L}}         % Likelihood
\newcommand{\Evid}{\mathcal{Z}}         % Evidence
\newcommand{\Post}{\mathcal{P}}         % Posterior

% Abbreviations for references
\newcommand{\CK}{C\&K}                  % Congdon & Keeton (2018)
\newcommand{\NB}{N\&B}                  % Narayan & Bartelmann (1997)
\newcommand{\Night}{Nightingale+18}     % Nightingale, Dye & Massey (2018)

% --- Headers ---
\pagestyle{fancy}
\fancyhf{}
\fancyhead[L]{\small\textit{Learning to Autolens}}
\fancyhead[R]{\small\thepage}
\renewcommand{\headrulewidth}{0.4pt}

% --- Theorem environments ---
\theoremstyle{definition}
\newtheorem{definition}{Definition}[section]
\newtheorem{example}{Example}[section]

% ============================================================

% --- Packages ---
\usepackage[utf8]{inputenc}
\usepackage[T1]{fontenc}
\usepackage{amsmath, amssymb, amsthm}
\usepackage{physics}       % \vb, \grad, \div, \curl, \dd, etc.
\usepackage{mathtools}     % \coloneqq, \prescript
\usepackage{bm}            % \bm for bold math
\usepackage{graphicx}
\usepackage{float}
\usepackage{booktabs}      % Professional tables
\usepackage{hyperref}
\usepackage{cleveref}      % \cref for smart cross-refs
\usepackage{xcolor}
\usepackage[margin=1in]{geometry}
\usepackage{enumitem}      % Custom list formatting
\usepackage{fancyhdr}      % Headers and footers
\usepackage{tcolorbox}     % Colored boxes for key results
\usepackage{listings}      % Code listings

% --- Color scheme ---
\definecolor{codeblue}{RGB}{31, 119, 180}
\definecolor{codegray}{RGB}{128, 128, 128}
\definecolor{keyresult}{RGB}{39, 123, 69}
\definecolor{exercise}{RGB}{178, 102, 0}

% --- tcolorbox environments ---
\newtcolorbox{keyresult}[1][]{
  colback=keyresult!5, colframe=keyresult!70,
  fonttitle=\bfseries, title={Key Result}, #1
}

\newtcolorbox{pythonbox}[1][]{
  colback=codeblue!5, colframe=codeblue!50,
  fonttitle=\bfseries, title={PyAutoLens Implementation}, #1
}

\newtcolorbox{exercisebox}[1][]{
  colback=exercise!5, colframe=exercise!50,
  fonttitle=\bfseries, title={Exercise}, #1
}

% --- Code listing style ---
\lstset{
  language=Python,
  basicstyle=\ttfamily\small,
  keywordstyle=\color{codeblue}\bfseries,
  commentstyle=\color{codegray}\itshape,
  stringstyle=\color{keyresult},
  breaklines=true,
  frame=single,
  framerule=0.5pt,
  rulecolor=\color{codegray!50},
}

% --- Physics macros ---
% Vectors (bold upright)
\newcommand{\vtheta}{\bm{\theta}}       % Image-plane position
\newcommand{\vbeta}{\bm{\beta}}         % Source-plane position
\newcommand{\valpha}{\bm{\alpha}}       % Deflection angle
\newcommand{\vx}{\bm{x}}               % Generic position

% Lensing quantities
\newcommand{\thetaE}{\theta_{\rm E}}    % Einstein radius
\newcommand{\SigmaCr}{\Sigma_{\rm cr}}  % Critical surface density
\newcommand{\Dd}{D_{\rm d}}             % Angular diameter distance to lens
\newcommand{\Ds}{D_{\rm s}}             % Angular diameter distance to source
\newcommand{\Dds}{D_{\rm ds}}           % Angular diameter distance lens-source

% Statistical quantities
\newcommand{\chisq}{\chi^2}             % Chi-squared
\newcommand{\chisqred}{\chi^2_{\rm red}} % Reduced chi-squared
\newcommand{\Like}{\mathcal{L}}         % Likelihood
\newcommand{\Evid}{\mathcal{Z}}         % Evidence
\newcommand{\Post}{\mathcal{P}}         % Posterior

% Abbreviations for references
\newcommand{\CK}{C\&K}                  % Congdon & Keeton (2018)
\newcommand{\NB}{N\&B}                  % Narayan & Bartelmann (1997)
\newcommand{\Night}{Nightingale+18}     % Nightingale, Dye & Massey (2018)

% --- Headers ---
\pagestyle{fancy}
\fancyhf{}
\fancyhead[L]{\small\textit{Learning to Autolens}}
\fancyhead[R]{\small\thepage}
\renewcommand{\headrulewidth}{0.4pt}

% --- Theorem environments ---
\theoremstyle{definition}
\newtheorem{definition}{Definition}[section]
\newtheorem{example}{Example}[section]

% ============================================================

% --- Packages ---
\usepackage[utf8]{inputenc}
\usepackage[T1]{fontenc}
\usepackage{amsmath, amssymb, amsthm}
\usepackage{physics}       % \vb, \grad, \div, \curl, \dd, etc.
\usepackage{mathtools}     % \coloneqq, \prescript
\usepackage{bm}            % \bm for bold math
\usepackage{graphicx}
\usepackage{float}
\usepackage{booktabs}      % Professional tables
\usepackage{hyperref}
\usepackage{cleveref}      % \cref for smart cross-refs
\usepackage{xcolor}
\usepackage[margin=1in]{geometry}
\usepackage{enumitem}      % Custom list formatting
\usepackage{fancyhdr}      % Headers and footers
\usepackage{tcolorbox}     % Colored boxes for key results
\usepackage{listings}      % Code listings

% --- Color scheme ---
\definecolor{codeblue}{RGB}{31, 119, 180}
\definecolor{codegray}{RGB}{128, 128, 128}
\definecolor{keyresult}{RGB}{39, 123, 69}
\definecolor{exercise}{RGB}{178, 102, 0}

% --- tcolorbox environments ---
\newtcolorbox{keyresult}[1][]{
  colback=keyresult!5, colframe=keyresult!70,
  fonttitle=\bfseries, title={Key Result}, #1
}

\newtcolorbox{pythonbox}[1][]{
  colback=codeblue!5, colframe=codeblue!50,
  fonttitle=\bfseries, title={PyAutoLens Implementation}, #1
}

\newtcolorbox{exercisebox}[1][]{
  colback=exercise!5, colframe=exercise!50,
  fonttitle=\bfseries, title={Exercise}, #1
}

% --- Code listing style ---
\lstset{
  language=Python,
  basicstyle=\ttfamily\small,
  keywordstyle=\color{codeblue}\bfseries,
  commentstyle=\color{codegray}\itshape,
  stringstyle=\color{keyresult},
  breaklines=true,
  frame=single,
  framerule=0.5pt,
  rulecolor=\color{codegray!50},
}

% --- Physics macros ---
% Vectors (bold upright)
\newcommand{\vtheta}{\bm{\theta}}       % Image-plane position
\newcommand{\vbeta}{\bm{\beta}}         % Source-plane position
\newcommand{\valpha}{\bm{\alpha}}       % Deflection angle
\newcommand{\vx}{\bm{x}}               % Generic position

% Lensing quantities
\newcommand{\thetaE}{\theta_{\rm E}}    % Einstein radius
\newcommand{\SigmaCr}{\Sigma_{\rm cr}}  % Critical surface density
\newcommand{\Dd}{D_{\rm d}}             % Angular diameter distance to lens
\newcommand{\Ds}{D_{\rm s}}             % Angular diameter distance to source
\newcommand{\Dds}{D_{\rm ds}}           % Angular diameter distance lens-source

% Statistical quantities
\newcommand{\chisq}{\chi^2}             % Chi-squared
\newcommand{\chisqred}{\chi^2_{\rm red}} % Reduced chi-squared
\newcommand{\Like}{\mathcal{L}}         % Likelihood
\newcommand{\Evid}{\mathcal{Z}}         % Evidence
\newcommand{\Post}{\mathcal{P}}         % Posterior

% Abbreviations for references
\newcommand{\CK}{C\&K}                  % Congdon & Keeton (2018)
\newcommand{\NB}{N\&B}                  % Narayan & Bartelmann (1997)
\newcommand{\Night}{Nightingale+18}     % Nightingale, Dye & Massey (2018)

% --- Headers ---
\pagestyle{fancy}
\fancyhf{}
\fancyhead[L]{\small\textit{Learning to Autolens}}
\fancyhead[R]{\small\thepage}
\renewcommand{\headrulewidth}{0.4pt}

% --- Theorem environments ---
\theoremstyle{definition}
\newtheorem{definition}{Definition}[section]
\newtheorem{example}{Example}[section]


\fancyhead[C]{\small\textbf{Module 02: Simulating Lens Data}}

\title{\textbf{Module 02: Simulating Lens Data}\\[0.5em]
\large Theory Companion — Learning to Autolens}
\author{Rodrigo C\'ordova Rosado\\
\small Harvard-Smithsonian Center for Astrophysics\\
\small \texttt{rodrigo.cordova\_rosado@cfa.harvard.edu}}
\date{\today}

\begin{document}
\maketitle

\begin{abstract}
This document covers the image formation model: how a true lensed
surface brightness distribution is transformed into a CCD image through
PSF convolution and noise addition. Understanding this forward model is
essential for the inverse problem of lens modeling (Module~03).
References: \Night\ Sec.~3; Meneghetti (2021) Ch.~8.
\end{abstract}

\tableofcontents
\newpage

% ============================================================
\section{The Forward Model}
\label{sec:forward-model}
% ============================================================

The observed image $\mathbf{d}$ relates to the true surface brightness
$\mathbf{I}$ by:

\begin{keyresult}
\begin{equation}
\mathbf{d} = \mathrm{PSF} \ast \mathbf{I}(\bm{\eta}) + \mathbf{n}
\label{eq:forward-model}
\end{equation}
\end{keyresult}

\noindent where $\bm{\eta}$ are the model parameters (lens mass, source
light, etc.), $\ast$ denotes 2D convolution, and $\mathbf{n}$ is the
noise realization (\Night\ eq.~1).

The true lensed image is computed by ray-tracing (Module~01):
\begin{equation}
I(\vtheta) = I_s\bigl(\vbeta(\vtheta)\bigr) + I_\ell(\vtheta)
\label{eq:lensed-image}
\end{equation}
where $I_s$ is the source surface brightness evaluated at the
deflected position $\vbeta = \vtheta - \valpha(\vtheta)$, and $I_\ell$
is the lens galaxy's own light.

% ============================================================
\section{Point Spread Function}
\label{sec:psf}
% ============================================================

\subsection{Gaussian Approximation}

For a circular Gaussian PSF with standard deviation $\sigma$:
\begin{equation}
\mathrm{PSF}(x, y) = \frac{1}{2\pi\sigma^2}
\exp\!\left(-\frac{x^2 + y^2}{2\sigma^2}\right)
\label{eq:gaussian-psf}
\end{equation}

The PSF is normalized: $\iint \mathrm{PSF}\,dx\,dy = 1$. The
\textbf{FWHM} (full width at half maximum) is:
\begin{equation}
\mathrm{FWHM} = 2\sqrt{2\ln 2}\,\sigma \approx 2.355\,\sigma
\label{eq:fwhm}
\end{equation}

\subsection{Convolution via Fourier Transform}

Convolution in real space corresponds to multiplication in Fourier space
(the convolution theorem):
\begin{equation}
\mathcal{F}[\mathrm{PSF} \ast I] = \mathcal{F}[\mathrm{PSF}] \cdot \mathcal{F}[I]
\label{eq:convolution-theorem}
\end{equation}

For a Gaussian PSF, the Fourier transform is also Gaussian:
$\mathcal{F}[\mathrm{Gauss}(\sigma_x)] = \mathrm{Gauss}(\sigma_k)$ with
$\sigma_k = 1/(2\pi\sigma_x)$. A narrow PSF (sharp image) preserves
high spatial frequencies; a broad PSF (blurry image) suppresses them.

\subsection{Nyquist Sampling}

To properly sample the PSF without aliasing, the pixel scale $\Delta\theta$
must satisfy:
\begin{equation}
\Delta\theta \leq \frac{\mathrm{FWHM}}{2}
\label{eq:nyquist}
\end{equation}

\subsection{Instrument Parameters}

\begin{table}[H]
\centering
\caption{Typical parameters for common instruments.}
\begin{tabular}{lcccc}
\toprule
Instrument & Pixel scale & PSF FWHM & Exposure & Sky level \\
\midrule
HST ACS/WFC3 & $0.05''$/pix & $0.07''$--$0.10''$ & 2000\,s & 0.01 \\
Euclid VIS   & $0.10''$/pix & $0.16''$           & 2260\,s & 0.04 \\
Keck AO      & $0.01''$--$0.04''$/pix & $0.05''$--$0.08''$ & 600\,s & 0.5 \\
Seeing       & $0.10''$--$0.25''$/pix & $0.5''$--$1.5''$ & 900\,s & 1.0 \\
\bottomrule
\end{tabular}
\label{tab:instruments}
\end{table}

% ============================================================
\section{Noise Model}
\label{sec:noise}
% ============================================================

\subsection{CCD Noise Components}

The total noise variance per pixel is:
\begin{equation}
\sigma_{\rm total}^2 = F_{\rm src}\,t\,\Omega_{\rm pix}
+ F_{\rm sky}\,t\,\Omega_{\rm pix} + \sigma_{\rm read}^2
\label{eq:noise-model}
\end{equation}
where $F_{\rm src}$ and $F_{\rm sky}$ are flux rates
(e$^-$/s/arcsec$^2$), $t$ is the exposure time, $\Omega_{\rm pix}$ is
the pixel solid angle, and $\sigma_{\rm read}$ is the read noise
(Howell 2006, Ch.~4).

\subsection{Signal-to-Noise Ratio}

The SNR per pixel is:
\begin{equation}
\mathrm{SNR} = \frac{F_{\rm src}\,t\,\Omega_{\rm pix}}{\sigma_{\rm total}}
\label{eq:snr}
\end{equation}

In the \textbf{background-limited} regime ($F_{\rm sky} \gg F_{\rm src}$),
$\mathrm{SNR} \propto \sqrt{t}$. In the \textbf{source-limited} regime,
$\mathrm{SNR} \propto \sqrt{t}$ as well (Poisson statistics).

% ============================================================
\section{Over-Sampling}
\label{sec:oversampling}
% ============================================================

A single pixel integrates the surface brightness over its area. For
steep profiles or thin arcs, evaluating $I(\vtheta)$ only at the pixel
center introduces bias. \textbf{Over-sampling} subdivides each pixel
into $N \times N$ sub-pixels, evaluates the profile at each, and
averages:
\begin{equation}
I_{\rm pixel} = \frac{1}{N^2} \sum_{k=1}^{N^2} I(\vtheta_k)
\label{eq:oversampling}
\end{equation}

PyAutoLens uses \textbf{adaptive radial over-sampling}: high $N$ near
the center (where profiles are steep), low $N$ in the outskirts.

\begin{pythonbox}
\begin{lstlisting}
over_sample = al.util.over_sample.over_sample_size_via_radial_bins_from(
    grid=grid,
    sub_size_list=[32, 8, 2],  # 32x32 near center, 2x2 at edge
    radial_list=[0.3, 0.6],   # Transition radii (arcsec)
    centre_list=[(0.0, 0.0)],
)
\end{lstlisting}
\end{pythonbox}

% ============================================================
\section*{References}
% ============================================================

\begin{itemize}[nosep, leftmargin=*]
  \item Howell, S.B.\ (2006), \emph{Handbook of CCD Astronomy}, Cambridge Univ.\ Press
  \item Meneghetti, M.\ (2021), \emph{Introduction to Gravitational Lensing}, Springer
  \item Nightingale, J.W., Dye, S.\ \& Massey, R.J.\ (2018), MNRAS, 478, 4738
\end{itemize}

\end{document}
